% Slide 13 — Portfólio de políticas
\begin{frame}{Portfólio de políticas candidatas}
  \centering
  \resizebox{\linewidth}{!}{%
  \begin{tabular}{@{}llll@{}}
    \toprule
    \textbf{Clássicas} & \textbf{Param. por meta-heurística} & \textbf{Aprendidas por reforço} & \textbf{Híbridas GA-RL} \\
    \midrule
    EOQ & Política parametrizada por GA & DQN & GA-DQN \\
    $(s,S)$ & Política parametrizada por SA & PPO & GA-PPO \\
    Jornaleiro & Política parametrizada por PSO; Política parametrizada por DE & SARSA & \\
    \bottomrule
  \end{tabular}}
  \vspace{1em}

  \begin{alertblock}{Cuidado conceitual}
    \small GA, SA, PSO e DE \highlight{não são políticas autônomas}: são
    métodos de busca usados para parametrizar a regra $(s,Q)$. Quem decide é
    sempre a regra de reposição; a meta-heurística apenas a calibra.
  \end{alertblock}
  \note{
    O portfólio reúne quatro famílias. As políticas clássicas são regras
    analíticas: EOQ, (s,S) e o Jornaleiro. As meta-heurísticas -- GA, SA, PSO
    e DE -- não são políticas no sentido decisório: elas são algoritmos de
    busca usados para encontrar os melhores parâmetros (s, Q) de uma regra
    de reposição, avaliada sempre pelo mesmo simulador. As políticas
    aprendidas por reforço -- DQN, PPO e SARSA -- modelam a reposição como
    um processo de decisão sequencial. E as híbridas GA-DQN e GA-PPO usam o
    GA para inicializar o treinamento do agente de RL, funcionando como
    especialistas do portfólio, não como a contribuição central do
    trabalho.
  }
\end{frame}
